% Compendium of the verified formulas: General Sinus, General Numbers, Cosmic Mass
% (book "The Cosmic Mass, Volume 1", and the three papers). 10 November 2026.
\documentclass[11pt]{article}
\usepackage[a4paper,margin=2.4cm]{geometry}
\usepackage{amsmath,amssymb,mathtools}
\usepackage{booktabs,array,tabularx}
\usepackage[hidelinks]{hyperref}
\usepackage{xcolor}

\newcommand{\pb}{\overline{\phi}}
\newcommand{\Gp}{\mathcal{G}}
\newcommand{\G}{\lvert\Gp\rvert}
\newcommand{\uE}{\boldsymbol{\mathcal{E}}}
\newcommand{\uB}{\boldsymbol{\mathcal{B}}}
\newcommand{\vv}{\mathbf{v}}
\newcommand{\vg}{\mathbf{g}}
\newcommand{\ve}{\mathbf{e}}
\newcommand{\vE}{\mathbf{E}}
\newcommand{\vB}{\mathbf{B}}
\renewcommand{\Re}{\operatorname{Re}}
\DeclareMathOperator{\Com}{Com}
\DeclareMathOperator{\Sh}{Sh}
\DeclareMathOperator{\Ht}{Ht}
\newcommand{\src}[1]{\hfill{\small\textcolor{gray}{[#1]}}}

\title{Verified formulas of the General Sinus, General Numbers\\ and Cosmic Mass theory}
\author{Work preparation package for the virtual universe}
\date{10 November 2026}

\begin{document}
\maketitle

\noindent Every formula below was checked by a script of \texttt{02-verification-scripts} or by
the 50-digit audit of this package. The tag in square brackets names the source:

\begin{itemize}
\item \textbf{GS:} General Sinus, revision 3, and the book's Appendix A;
\item \textbf{GN:} the General Numbers paper and the book's Chapter 4;
\item \textbf{CM:} the Cosmic Mass paper and the book's Chapters 6 and 9.
\end{itemize}

\noindent Formulas found false are listed only in Section~\ref{sec:corrections}, with
their corrections.

%=====================================================================
\section{General Sinus: trigonometry of an oblique projection}
%=====================================================================

\paragraph{Definitions.} The point $Z$ on the second side of the angle $\alpha$ is
projected onto the first side along a line at the angle $\phi$ to $ZO$:
\[
\sin(\alpha,\phi) = \frac{\sin\alpha}{\sin(\alpha+\phi)},\quad
\cos(\alpha,\phi) = \frac{\sin\phi}{\sin(\alpha+\phi)},\quad
\tan(\alpha,\phi) = \frac{\sin\alpha}{\sin\phi},
\]
\[
\csc(\alpha,\phi) = \frac{\sin(\alpha+\phi)}{\sin\alpha},\quad
\sec(\alpha,\phi) = \frac{\sin(\alpha+\phi)}{\sin\phi},\quad
\cot(\alpha,\phi) = \frac{\sin\phi}{\sin\alpha}.
\src{GS}
\]
For $\alpha+\phi = \pi/2$ the classical functions return.

\paragraph{Coordinate reading and Pythagoras.} With $\pb = \alpha+\phi$ and
$\vec u = (\cos\pb,\sin\pb)$,
\begin{gather*}
(\cos\alpha,\sin\alpha) = \cos(\alpha,\phi)\,\vec e_1 + \sin(\alpha,\phi)\,\vec u ,\\
\cos(\alpha,\phi)^2 + \sin(\alpha,\phi)^2 + 2\cos(\alpha,\phi)\sin(\alpha,\phi)\cos(\alpha+\phi) = 1 .
\end{gather*}\src{GS}

\paragraph{Exchange, reciprocals, signs.}
\[
\sin(\alpha,\phi) = \cos(\phi,\alpha),\quad \tan(\alpha,\phi) = \cot(\phi,\alpha),\quad
\csc(\alpha,\phi) = \sec(\phi,\alpha);
\]
\[
\csc(\alpha,\phi) = \sin(\alpha+\phi,-\phi),\quad
\sec(\alpha,\phi) = \tan(\alpha+\phi,\phi),\quad
\cot(\alpha,\phi) = \cos(\alpha-\phi,\phi);
\]
\begin{gather*}
\sin(-\alpha,-\phi) = \sin(\alpha,\phi),\quad \cos(-\alpha,-\phi) = \cos(\alpha,\phi),\quad
\sin(-\alpha,\phi) = \sin(\alpha,-\phi),\\ \tan(i\alpha,j\phi) = ij\tan(\alpha,\phi)\quad (i,j = \pm1);\\
\sin(\alpha+p\pi,\phi+q\pi) = (-1)^q\sin(\alpha,\phi),\quad
\cos(\alpha+p\pi,\phi+q\pi) = (-1)^p\cos(\alpha,\phi),\\
\tan(\alpha+p\pi,\phi+q\pi) = (-1)^{p+q}\tan(\alpha,\phi).
\end{gather*}\src{GS}

\paragraph{Special values.}
$\sin(\alpha,\tfrac\pi2) = \tan\alpha$, $\cos(\alpha,\tfrac\pi2) = \sec\alpha$,
$\tan(\alpha,\tfrac\pi2) = \sin\alpha$; $\sin(\theta,\theta) = \cos(\theta,\theta) = \tfrac{1}{2\cos\theta}$,
$\tan(\theta,\theta) = 1$; all six equal $1$ at $\theta = \pi/3$; $\sin(\tfrac\pi5,\tfrac\pi5) = \tau^{-1}$
with $\tau = (1+\sqrt5)/2$. \src{GS}

\paragraph{Ptolemy's identity.}
$\sin(x+y)\sin(y+z) = \sin x\sin z + \sin y\sin(x+y+z)$. Every addition formula below reduces to it. \src{GS}

\paragraph{Addition with the projection angle fixed (moving frame).}
\begin{align*}
\sin(\alpha+\beta,\phi) &= \sin(\alpha,\phi) + \cos(\alpha,\phi)\sin(\beta,\phi+\alpha)\\
 &= \sin(\alpha,\phi) + \cos(\alpha,\phi)\sin(\beta,\phi)\cos(\alpha,\phi+\beta),\\
\cos(\alpha+\beta,\phi) &= \cos(\alpha,\phi)\cos(\beta,\phi) + \sin(\alpha,\phi)\sin(\beta,\phi)\cos(\alpha+\beta,\phi)\\
 &= \frac{\cos(\alpha,\phi)\cos(\beta,\phi)}{1-\sin(\alpha,\phi)\sin(\beta,\phi)}\\
 &= \cos(\alpha,\phi)\cos(\beta,\phi) + \cos(\alpha,\phi)\sin(\beta,\phi)\sin(\alpha,\phi+\beta),\\
\sec(\alpha+\beta,\phi) &= \sec(\alpha,\phi)\sec(\beta,\phi) - \tan(\alpha,\phi)\tan(\beta,\phi),\\
\tan(\alpha+\beta,\phi) &= \tan(\alpha,\phi)\cos\beta + \tan(\beta,\phi)\cos\alpha .
\end{align*}\src{GS}

\paragraph{Addition in the projection angle.}
$\sin(\alpha,\phi+\theta) = \sin(\alpha,\phi)\sin(\alpha+\phi,\theta)$,\quad
$\cos(\alpha,\phi+\theta) = \cos(\alpha,\phi) + \sin(\alpha,\phi)\cos(\alpha+\phi,\theta)$. \src{GS}

\paragraph{Double and half angles.}
\begin{gather*}
\sin(2\alpha,\phi) = \sin(\alpha,\phi)\bigl(1+\cos(2\alpha,\phi)\bigr),\quad
\cos(2\alpha,\phi) = \frac{\cos(\alpha,\phi)^2}{1-\sin(\alpha,\phi)^2},\\
\sec(2\alpha,\phi) = \sec(\alpha,\phi)^2 - \tan(\alpha,\phi)^2,\\
\tan(2\alpha,\phi) = 2\tan(\alpha,\phi)\cos\alpha,\quad
\sin(2\alpha,\phi)-\sin(\alpha,\phi) = \sin(\alpha,\phi)\cos(2\alpha,\phi),\\
\cos(\alpha,\phi) = \frac{\cos(\frac\alpha2,\phi)^2}{1-\sin(\frac\alpha2,\phi)^2},\quad
\sec(\alpha,\phi) = \sec(\tfrac\alpha2,\phi)^2 - \tan(\tfrac\alpha2,\phi)^2 .
\end{gather*}
\src{GS}

\paragraph{Fixed frame (Euler and De Moivre).} For fixed $\pb\in(0,\pi)$, write
$C_\theta = \sin(\pb-\theta)/\sin\pb$, $S_\theta = \sin\theta/\sin\pb$ and $u = e^{i\pb}$. Then
\begin{gather*}
e^{i\theta} = C_\theta + S_\theta u,\qquad u^2 = 2\cos\pb\,u - 1,\\
C_{\alpha+\beta} = C_\alpha C_\beta - S_\alpha S_\beta,\qquad
S_{\alpha+\beta} = C_\alpha S_\beta + S_\alpha C_\beta + 2\cos\pb\,S_\alpha S_\beta,\\
(C_\theta + S_\theta u)^n = C_{n\theta} + S_{n\theta}u,\\
u^n = -U_{n-2}(\cos\pb) + U_{n-1}(\cos\pb)\,u\quad(n\ge 2).
\end{gather*}
\src{GS}

\paragraph{Partial derivatives (corrected second order).}
\begin{gather*}
\partial_x\sin(x,y) = \frac{\cos(x,y)}{\sin(x+y)},\quad \partial_y\sin(x,y) = -\sin(x,y)\cot(x+y),\\
\partial_x\cos(x,y) = -\cos(x,y)\cot(x+y),\quad
\partial_y\cos(x,y) = \frac{\sin(x,y)}{\sin(x+y)},\\
\partial_x\tan(x,y) = \frac{\cos x}{\sin y} = \tan(\tfrac\pi2\pm x,y),\quad \partial_y\tan(x,y) = -\tan(x,y)\cot y,\\
\partial_x^2\sin = -2\cot(x+y)\,\partial_x\sin,\quad \partial_y^2\sin = \sin(x,y)\bigl(1+2\cot^2(x+y)\bigr),\\
\partial_x^2\cos = \cos(x,y)\bigl(1+2\cot^2(x+y)\bigr),\quad
\partial_y^2\cos = -2\cot(x+y)\,\partial_y\cos,\\ \partial_x^2\tan = -\tan(x,y),\quad
\partial_y^2\tan = \tan(x,y)\bigl(1+2\cot^2 y\bigr).
\end{gather*}
\src{GS}

\paragraph{Derivative with the frame fixed.} $Df = \partial_x f - \partial_y f$ (path $x+t$, $y-t$). For every integer $n\ge0$:
\begin{gather*}
D^n\sin(x,y) = \sin\bigl(x+\tfrac{n\pi}{2},\,y-\tfrac{n\pi}{2}\bigr),\qquad
D^n\cos(x,y) = \cos\bigl(x+\tfrac{n\pi}{2},\,y-\tfrac{n\pi}{2}\bigr),\qquad D^2 f = -f,\\
D^n\sin\bigl(x-\tfrac{n\pi}{2},\,y+\tfrac{n\pi}{2}\bigr) = \sin(x,y)\ \text{(primitives)},\\
D^n\tan(x,y) = \sin(x+y)\Bigl(-\frac{d}{dy}\Bigr)^{n-1}\csc^2 y\quad(n\ge1).
\end{gather*}
\src{GS}

\paragraph{Products for every $N\ge1$.}
\begin{gather*}
\sin(x,y) = \prod_{k=0}^{N-1}\sin\Bigl(\tfrac{x+\pi k}{N},\tfrac{y}{N}\Bigr),\quad
\cos(x,y) = \prod_{k=0}^{N-1}\cos\Bigl(\tfrac{x}{N},\tfrac{y+\pi k}{N}\Bigr),\\
\tan(x,y) = \prod_{k=0}^{N-1}\tan\Bigl(\tfrac{x+\pi k}{N},\tfrac{y+\pi k}{N}\Bigr).
\end{gather*}\src{GS}

\paragraph{General Sinus Transform.} Take $P_{N,k}(z) = u_k z^2 - 2Nv_kz - 2N^2u_k$ with exact $u_k = \sin(\pi k/N)$ and $v_k = \cos(\pi k/N)$, and set $S_N = \prod_k P_{N,k}(x)/P_{N,k}(x+y)$ ($C_N$ and $T_N$ likewise). With $\rho(\varepsilon) = \sqrt{1+\varepsilon^4/4}$ and $\delta(\varepsilon) = \arctan\frac{\varepsilon}{1-\varepsilon^2/2}$, for $N > \max(|x|,|y|,|x+y|)/\sqrt2$:
\begin{gather*}
S_N(x,y) = \Bigl(\frac{\rho(x/N)}{\rho((x+y)/N)}\Bigr)^N\frac{\sin(N\delta(x/N))}{\sin(N\delta((x+y)/N))},\\
S_N = \sin(x,y) + \frac{x^3\cos x\sin(x+y) - (x+y)^3\sin x\cos(x+y)}{6N^2\sin^2(x+y)} + O(N^{-3}).
\end{gather*}\src{GS}

%=====================================================================
\section{General Numbers: the plane in a leaning frame}
%=====================================================================

\paragraph{Axiom and the two angles of one length.}
$\Gp = \G\,\vec u$, $\vec u = (\cos\pb,\sin\pb)$, $\G = \lvert\cos\pb+\sin\pb\rvert = \sqrt2\,\lvert\sin(\pb+\tfrac\pi4)\rvert$.
For $1<A<\sqrt2$: $\pb_1 = \tfrac12\arcsin(A^2-1)$, $\pb_2 = \tfrac\pi2-\pb_1$ ($A = 1.16043$: $10.14^\circ$ and $79.86^\circ$). \src{GN}

\paragraph{Rule of the Square and the product.} For $z = a+b\Gp$ and $w = c+d\Gp$:
\[
\Gp^2 = 2\G\cos\pb\,\Gp - \G^2,\qquad
zw = (ac - \G^2bd) + (ad+bc+2\G\cos\pb\,bd)\,\Gp .
\src{GN}
\]

\paragraph{Readings.}
$\Re z = a$, $\Com z = b\Gp$, $\Sh z = a + b\G\cos\pb$, $\Ht z = b\G\sin\pb$, and
$\lvert z\rvert^2 = a^2+2ab\G\cos\pb+b^2\G^2 = (\Sh z)^2+(\Ht z)^2$. \src{GN}

\paragraph{Conjugate, skew flip, division.}
$\bar z = (a+2b\G\cos\pb) - b\Gp$, $z+\bar z = 2\Sh z$, $\Re\bar z = 2\Sh z - \Re z$,
$z\bar z = \lvert z\rvert^2$, $1/z = \bar z/\lvert z\rvert^2$. The skew flip is $a - b\Gp$, and
$(a+b\Gp)(a-b\Gp) = a^2+b^2\G^2 - 2b^2\G\cos\pb\,\Gp$. The norm is multiplicative. \src{GN}

\paragraph{Turning Unit and the bridge to $\mathbb C$.}
$\mathcal J = -\cot\pb + \Gp/(\G\sin\pb)$, $\mathcal J^2 = -1$, $z = \Sh z + \Ht z\,\mathcal J$; the map
$z\mapsto\Sh z + i\Ht z$ is an isomorphism of fields onto $\mathbb C$. \src{GN}

\paragraph{Laws of the parts.}
$\Re(zw) = \Re z\Re w - \G^2bd$, $\Com(zw) = \Sh(z)\Com(w)+\Sh(w)\Com(z)$,
$\Sh(zw) = \Sh z\Sh w - \Ht z\Ht w$, $\Ht(zw) = \Sh z\Ht w + \Sh w\Ht z$,
$\lvert\Re z\rvert \le \lvert z\rvert/\sin\pb$ (the bound is reached). \src{GN}

\paragraph{Turning, powers, roots.}
$e^{\theta\mathcal J} = C_\theta + S_\theta\vec u$; $(C_\theta+S_\theta\vec u)^n = C_{n\theta}+S_{n\theta}\vec u$;
$\Gp^n = -\G^n\frac{\sin(n-1)\pb}{\sin\pb} + \G^{n-1}\frac{\sin n\pb}{\sin\pb}\Gp$;
the powers lie on the spiral $\rho = \G^{\theta/\pb}$; the $n$-th roots are $\lvert z\rvert^{1/n}(C_{\theta_k}+S_{\theta_k}\vec u)$
with $\theta_k = (\alpha+2\pi k)/n$. \src{GN}

\paragraph{Cauchy--Riemann in a leaning frame.} For $f = U+V\Gp$ holomorphic:
$\partial_bU = -\G^2\partial_aV$ and $\partial_bV = \partial_aU + 2\G\cos\pb\,\partial_aV$. \src{GN}

\paragraph{Generalized complex numbers.} $\Gp^2 = p+q\Gp$ with $p = -\G^2$ and $q = 2\G\cos\pb$, so $q^2+4p<0$ (the elliptic case).
Conversely, every elliptic pair $(p,q)$ is a leaning frame with length $\sqrt{-p}$ and angle $\arccos(q/2\sqrt{-p})$. \src{GN}

%=====================================================================
\section{Cosmic Mass}
%=====================================================================

\paragraph{The cosmic mass and its unit.}
$\mu = m + q\,\Gp$, with $\G = s = \sqrt{K/G} = 1.16043\times10^{10}$ kg/C. The drawing length is $A = s/u$, and $\lvert\cos\pb+\sin\pb\rvert = A$; with $u = 10^{10}$ kg/C, $A = 1.16043$ and $\pb = 79.86^\circ$ (or $10.14^\circ$). In Planck or Stoney units $A = 1$ and $\pb = 90^\circ$. \src{CM}

\paragraph{Cosmic Interaction Law and its reading.}
\begin{gather*}
\vec F_G = \frac{G\,\mu_a\mu_b}{r^2}\,\vec e_r,\qquad
F = \frac{G\,\Re(\mu_a\mu_b)}{r^2} = \frac{G m_am_b - K q_aq_b}{r^2}\ \text{ at every angle},\\
\Re\bigl[(a+b\Gp)(c+d\Gp)\bigr] = ac - s^2bd .
\end{gather*}\src{CM}

\paragraph{Full product (both coordinates).}
\[
\mu_a\mu_b = (m_am_b - s^2q_aq_b) + \bigl(m_aq_b + q_am_b + 2s\cos\pb\,q_aq_b\bigr)\Gp .
\]
The Composer coordinate holds the Coupling term $m_aq_b+q_am_b$ and the charge share $2s\cos\pb\,q_aq_b$.
It is not read as a force: see the book's Chapter 10 and Section~\ref{sec:corrections}. \src{CM}

\paragraph{Stoney mass.} $m_S = e\,s \approx 1.85921\times10^{-9}$ kg $=\sqrt{\alpha_{\mathrm{fs}}}\,m_P$. \src{CM}

\paragraph{Cosmic Field Law and law of motion.} Fields with General Number values,
$\uE = \vg + \ve\,\Gp$ and $\uB = \vB_g + \mathbf b\,\Gp$; sources
$\rho_\mu = \rho_m+\rho_q\Gp$ and $\mathbf J_\mu = \mathbf J_m+\mathbf J_q\Gp$:
\begin{gather*}
\nabla\cdot\uE = -4\pi G\rho_\mu,\quad \nabla\cdot\uB = 0,\quad \nabla\times\uE = -\partial_t\uB,\quad
\nabla\times\uB = -\frac{4\pi G}{c^2}\mathbf J_\mu + \frac{1}{c^2}\partial_t\uE,\\
\Re(\mu)\frac{d\vv}{dt} = \Re\bigl[\mu(\uE+\vv\times\uB)\bigr]
\iff m\frac{d\vv}{dt} = m(\vg+\vv\times\vB_g) + q(\vE+\vv\times\vB).
\end{gather*}
Each field equation splits into its two coordinates. The Composer set is Maxwell's, with $\vE = -s^2\ve$, $\vB = -s^2\mathbf b$ and $\mu_0 = 4\pi K/c^2$. The field of a source at rest is $\uE = -G\mu_s\vec e_r/r^2$. \src{CM}

\paragraph{Kept quantities.}
$\partial_t\rho_\mu + \nabla\cdot\mathbf J_\mu = 0$. Both fields obey the wave equation with speed $c$.
\[
u = -\frac{g^2+c^2B_g^2}{8\pi G} + \frac{E^2+c^2B^2}{8\pi K},\qquad
\mathbf S = -\frac{c^2}{4\pi G}\Re(\uE\times\uB),\qquad \partial_tu + \nabla\cdot\mathbf S = -\Re(\uE\cdot\mathbf J_\mu).
\]
The momentum balance holds with the stress $T_{ij}$ of the paper. The action is
\[
S = \int dt\,\Bigl\{\sum\bigl[\tfrac12\Re(\mu)v^2 - \Re\bigl(\mu(\Phi-\vv\cdot\boldsymbol{\mathcal A})\bigr)\bigr]
- \frac{1}{8\pi G}\int\Re\bigl[\uE\cdot\uE - c^2\uB\cdot\uB\bigr]\,dV\Bigr\}.
\]\src{CM}

\paragraph{Where the law fails (stated in the paper).}
\begin{itemize}
\item Frame dragging on LAGEOS comes out at $1/4$ of the measured value.
\item There is no light bending and no Shapiro delay.
\item Gravity waves carry negative energy.
\item There is no conversion of rest mass into energy.
\item The law is valid for slow bodies only.
\end{itemize}
\src{CM}

%=====================================================================
\section{Corrections and exclusions}\label{sec:corrections}
%=====================================================================

{\small
\begin{tabularx}{\textwidth}{@{}>{\raggedright\arraybackslash}X >{\raggedright\arraybackslash}p{0.12\textwidth} >{\raggedright\arraybackslash}X@{}}
\toprule
\textbf{statement} & \textbf{status} & \textbf{correct form / note}\\
\midrule
$\tan(2\alpha,\phi)-\sec(2\alpha,\phi) = \sin(\alpha,\phi)$ & false & deleted\\
$\tan(\alpha,\phi) = \sin(\tfrac\alpha2,\phi)+\sec(\alpha,\phi)$ & false & deleted\\
$\partial_y^2\sin = \sin\,(2\cot^2(x+y)-1)$ & false & $\sin\,(1+2\cot^2(x+y))$\\
$\partial_x^2\cos = \cos\,(2\cot^2(x+y)-1)$ & false & $\cos\,(1+2\cot^2(x+y))$\\
$\partial_y^2\tan = \tan\,(2\cot^2(x+y)-1)$ & false & $\tan\,(1+2\cot^2y)$\\
$D^2\sin = +\sin$, $D^2\cos = +\cos$ & false & $D^2 f = -f$\\
$D^n$ and primitives by $\omega(z,n)$ & false for $n\equiv2,3$ & $D^nf(x,y) = f(x+\tfrac{n\pi}2,y-\tfrac{n\pi}2)$\\
$\tan^{(n)} = n!\,\tan'\tan(\tfrac\pi2+y,y)^{n-1}$ & true only $n\le2$ & $D^n\tan = \sin(x+y)(-\tfrac{d}{dy})^{n-1}\csc^2y$\\
rotation recurrence for $u_{n,k}$, $v_{n,k}$ ``exact'' & false & use exact sines; error accumulates otherwise\\
five digits for $n\ge20$ at every point & false with the recurrence & rate $c(x,y)/N^2$, proved (Section 1)\\
heat-equation analogy, $R^2 = 0.83$ & withdrawn & not evidence\\
Composer coordinate as a force along the line, full size & excluded by measurement & neutron would carry $9.0\times10^{-19}e$; measured $(-0.2\pm0.8)\times10^{-21}e$; share below $\approx0.002$ (book, Ch.~10)\\
charge share as a force along the line & excluded & hydrogen would bind at 5.71 eV, measured 13.598 eV (book, Ch.~10)\\
Standard Model (book, Ch.~7), Closure Law (Ch.~8) & not verified & kept out of the papers; Closure Law is a refitted empirical law\\
\bottomrule
\end{tabularx}}

\end{document}
